\section{KV-Cache}

\subsection{为什么需要 KV-Cache？}

\begin{importantbox}{KV-Cache 的核心动机}
在自回归生成中，每个新 token 需要与\textbf{所有历史 token} 计算 Attention。如果每步都重新计算所有 $K, V$，计算量为 $O(M \cdot d_k)$（其中 $M$ 不断增长）。KV-Cache 将历史 $K, V$ 向量缓存在 GPU 显存中，避免重复计算。
\end{importantbox}

\subsection{KV-Cache 的显存消耗}

每层每个 token 的 KV-Cache 大小：
$$
\text{Per-token KV} = 2 \times n_{\text{kv\_heads}} \times d_h \times \text{bytes}
$$

总 KV-Cache：
$$
\text{Total KV} = L \times M \times 2 \times n_{\text{kv\_heads}} \times d_h \times \text{bytes}
$$

其中 $L$ 是层数，$M$ 是序列长度。这就是为什么 GQA（减少 KV heads）对推理效率如此重要。

\subsection{本章小结}

KV-Cache 是自回归推理的核心优化，避免了重复计算 $K, V$。它的显存消耗随序列长度线性增长，是长上下文推理的主要瓶颈。

